\section{Literature Review}

\subsection{Academic Performance}

Academic performance is commonly used as an indicator of students' educational outcomes. It may be measured using examination scores, continuous assessment results, grade point averages, or other academic indicators \citep{determinants2023academic}.

\subsection{Factors Associated with Academic Performance}

Several factors may influence students' academic outcomes. These factors can be broadly classified into academic, personal, social, and demographic characteristics \citep{determinants2023academic}. Study habits, study skills, and attitudes towards learning have also been identified as important predictors of academic performance \citep{crede2008study}.

These factors may interact with one another and may operate through different aspects of students' academic and personal experiences. For example, time management, attendance, sleep, and self-esteem have been examined together as factors that may be associated with academic performance among university students \citep{masnan2025academic}. Considering these factors together provides a broader basis for understanding academic performance than focusing on a single characteristic in isolation.

\subsection{Theoretical Framework: Maslow's Hierarchy of Needs}

The range of personal and environmental factors associated with academic performance suggests that students' educational outcomes cannot always be considered independently of their broader needs and circumstances. Maslow's Hierarchy of Needs provides a theoretical perspective for considering how different human needs may influence motivation and behaviour \citep{freitas2011maslow}. The theory proposes five broad categories of needs: physiological needs, safety needs, love and belonging, esteem, and self-actualisation. In an educational context, these needs provide a useful framework for considering factors that may influence students' ability to engage effectively in learning \citep{freitas2011maslow}.

Physiological needs include basic requirements such as adequate food, rest, and physical wellbeing. Safety needs relate to security and freedom from conditions that may cause significant physical or psychological distress. Love and belonging involve relationships with family, peers, and other members of the learning community, while esteem relates to confidence, recognition, and a sense of competence. Self-actualisation represents the pursuit of personal growth and the achievement of one's potential.

The application of Maslow's framework to education suggests that students' academic experiences may be influenced by conditions extending beyond the classroom. These conditions may affect students' motivation, engagement, and ability to participate effectively in learning. Thus, the framework provides a useful perspective for considering the broader personal and environmental factors associated with academic performance \citep{freitas2011maslow}.

In this study, Maslow's Hierarchy of Needs is used as a theoretical lens rather than as a direct statistical model. It provides a broader perspective for understanding how personal and environmental conditions may relate to academic performance, while the statistical analysis will focus on measurable variables available in the study dataset.

\subsection{Study Time and Performance}

Study time represents the amount of time a student spends engaging with academic materials outside scheduled classes. Although study time is an important aspect of students' academic behaviour, research suggests that the amount of time spent studying does not necessarily translate directly into higher academic performance. The effectiveness and quality of study practices may be important considerations alongside the amount of time spent studying \citep{plant2005study}.

Study time should also be considered alongside the strategies students use during learning. The effectiveness of learning may depend not only on the amount of time devoted to academic activities but also on how that time is used. Research examining classroom attendance and learning strategies suggests that both behavioural and learning-related factors can contribute to differences in learning outcomes \citep{tetteh2018attendance}. This supports the consideration of study behaviour as a broader construct when examining factors associated with academic performance.

\subsection{Attendance and Performance}

Class attendance is another measurable factor that may be associated with academic performance. Research has examined the relationship between attendance and academic outcomes, with findings suggesting that attendance can be an important predictor of student performance \citep{arnold2018attendance}. The timing and pattern of student absences may also influence the relationship between absenteeism and academic achievement \citep{keppens2023absence}.

The relationship between attendance and academic performance may also depend on the learning context. Evidence from synchronous university courses indicates that class attendance can be associated with academic performance, suggesting that participation in scheduled learning activities may be relevant to students' academic outcomes \citep{ha2024attendance}. These findings further support the inclusion of attendance as a measurable factor when examining student academic performance.

\subsection{Statistical Modelling}

The discussion of study behaviour, attendance, and other factors associated with academic performance highlights the need for an approach that can examine several explanatory variables simultaneously. Statistical modelling provides such a framework by allowing relationships between a dependent variable and one or more explanatory variables to be examined systematically. In this study, multiple linear regression is considered as a suitable approach for modelling academic performance as a function of measurable student characteristics.

\vspace{0.5em}

\noindent A general multiple linear regression model can be expressed as

\[
Y = \beta_0 + \beta_1X_1 + \beta_2X_2 + \cdots + \beta_pX_p + \varepsilon,
\]

where \(Y\) represents academic performance, \(X_1,\ldots,X_p\) represent explanatory variables, \(\beta_0\) is the intercept, \(\beta_1,\ldots,\beta_p\) are regression coefficients, and \(\varepsilon\) is the error term. Previous research has demonstrated the use of multiple linear regression and principal component analysis for predicting students' academic performance \citep{yang2018predicting}.